% ============================================================
% §2 Preliminaries and Problem Setting — NEW section (OR convention)
% Formal definitions, notation, background extracted from method_framework
% ============================================================
\section{Preliminaries and Problem Setting}
\label{sec:preliminaries}

We first review diffusion models and fix notation, then formalize the
noise trajectory search problem that motivates our approach.

% -----------------------------------------------------------
\subsection{Diffusion Models}
\label{sec:diffusion-models}

Diffusion models generate samples by progressively denoising random noise into
data.
We briefly review the two most common formulations, as our framework applies to
both.

\subsubsection{The DDPM Formulation.}
The Denoising Diffusion Probabilistic Modeling (DDPM) formulation
\citep{ho2020denoising} couples two Markov processes: a forward process
($t = 1$ to $T$) that perturbs data with increasing Gaussian
noise, and a backward process ($t = T$ to $1$) that converts noise back
to data.
Given a data distribution $\mathbf{x}_0 \sim q_{\mathrm{data}}(\mathbf{x}_0)$,
the forward process progressively injects noise via
\begin{equation}
  q(\mathbf{x}_t \mid \mathbf{x}_{t-1})
  = \mathcal{N}\bigl(\mathbf{x}_t;\, \sqrt{1 - \beta_t}\,\mathbf{x}_{t-1},\,
    \beta_t \mathbf{I}\bigr),
  \label{eq:ddpm-forward}
\end{equation}
where $\beta_t \in (0,1)$ are fixed hyperparameters.
Setting $\alpha_t = 1 - \beta_t$ and
$\bar{\alpha}_t = \prod_{s=0}^{t}\alpha_s$, we can express the noisy sample at
any timestep directly as
\begin{equation}
  \mathbf{x}_t
  = \sqrt{\bar{\alpha}_t}\,\mathbf{x}_0
    + \sqrt{1 - \bar{\alpha}_t}\,\boldsymbol{\varepsilon},
  \qquad \boldsymbol{\varepsilon} \sim \mathcal{N}(\mathbf{0}, \mathbf{I}).
  \label{eq:ddpm-xt}
\end{equation}
A neural network
$\boldsymbol{\varepsilon}_\theta(\mathbf{x}_t, t)$ trained to predict the noise
parameterizes the backward process; samples emerge by iterating
\begin{equation}
  \mathbf{x}_{t-1}
  = \frac{1}{\sqrt{1 - \beta_t}}
    \biggl(\mathbf{x}_t - \frac{\beta_t}{\sqrt{1 - \bar{\alpha}_t}}\,
    \boldsymbol{\varepsilon}_\theta(\mathbf{x}_t, t)\biggr)
    + \sigma_t\,\boldsymbol{\varepsilon},
  \label{eq:ddpm-reverse}
\end{equation}
from $\mathbf{x}_T \sim \mathcal{N}(\mathbf{0}, \mathbf{I})$ down to
$\mathbf{x}_0$.

\subsubsection{The SGM Formulation.}
The Score-based Generative Model (SGM) formulation~\citep{song2019generative}
builds on the score function
$s(\mathbf{x}) = \nabla_\mathbf{x} \log q(\mathbf{x})$,
i.e., the gradient of the log-density.
Data are perturbed at different noise levels:
$\tilde{\mathbf{x}}_t = \mathbf{x}_0 + \sigma_t\boldsymbol{\varepsilon}$
for $0 < \sigma_1 < \sigma_2 < \cdots < \sigma_T$.
A neural network
$\mathbf{s}_\theta(\tilde{\mathbf{x}}_t, \sigma_t)$ approximates
the score at each noise level.
Annealed Langevin dynamics then generate samples: for each level
$t = T, T{-}1, \ldots, 1$, the sampler iterates
\begin{equation}
  \mathbf{x}_{(i)}
  = \mathbf{x}_{(i-1)}
    + \delta_i\,\nabla_\mathbf{x}\log q(\mathbf{x}_{(i-1)})
    + \sqrt{2\delta_i}\,\boldsymbol{\varepsilon}_i,
  \label{eq:sgm-langevin}
\end{equation}
starting from a prior distribution $\mathbf{x}_{(0)} \sim \pi$.

\begin{remark}
Both DDPM and SGM traverse $T$ discrete timesteps.
At each step, the transition from $\mathbf{x}_t$ to $\mathbf{x}_{t-1}$ depends
on a stochastic noise variable (the injected $\boldsymbol{\varepsilon}$ in
\eqref{eq:ddpm-reverse} or the Langevin noise in~\eqref{eq:sgm-langevin}).
This shared structure underpins the transition abstraction
in \cref{sec:framework}.
\end{remark}

% -----------------------------------------------------------
\subsection{Inference-Time Noise Search}
\label{sec:noise-search-problem}

\begin{definition}[Noise Trajectory Search Problem]
\label{def:nts}
Given a pretrained diffusion model with transition function
$F_\theta$, a conditioning signal $c$, a verifier $v(\cdot, c)$, and a total
compute budget of $B$ function evaluations (NFE), the \emph{noise trajectory
search problem} is to find a sequence of noise variables
$\{\varepsilon_t^*\}_{t=1}^T$ that maximizes the verifier score of the final
generated sample:
\begin{equation}
  \max_{\{\varepsilon_t\}_{t=1}^T}\;
  v\!\bigl(x_0(\varepsilon_1, \ldots, \varepsilon_T),\, c\bigr)
  \quad \text{subject to} \quad
  \text{total NFE} \leq B.
  \label{eq:nts-problem}
\end{equation}
\end{definition}

The core difficulty is sequential: the noise
$\varepsilon_t$ at timestep $t$ shapes the latent $x_{t-1}$, which in turn
determines which noises are effective at subsequent steps.
Because the budget $B$ must be respected globally, the problem further requires
allocating computation across timesteps.

\begin{definition}[Verifier]
\label{def:verifier}
A \emph{verifier} $v : \mathcal{X} \times \mathcal{C} \to \mathbb{R}$ scores a
generated sample $x_0$ against its conditioning signal $c$, returning a scalar
quality measure.
The verifier is fixed and not trained during search.
Examples include perceptual quality metrics, task-specific reward functions,
and domain-specific performance measures.
\end{definition}

\begin{definition}[Function Evaluation Budget]
\label{def:nfe}
A single \emph{function evaluation} (NFE) corresponds to one forward pass
through the denoising network.
The budget $B$ caps how many forward passes the sampler may
perform when generating a single sample.
Standard diffusion sampling consumes exactly $T$ NFE (one per timestep);
noise trajectory search spends the remaining $B - T$ NFE on local noise
refinement.
\end{definition}
